\section{Avancement et Architecture Technique}

\subsection{Structure et Interface de Base (thread.h)}
L'ensemble de l'interface de gestion de threads a été implémentée.

Tous les tests fournis passent à l'exception du test 71 (préemption, non implémentée).

\subsubsection*{Structure Fondamentale}
La structure \texttt{thread\_s} centralise toutes les informations :
\begin{itemize}
    \item Contexte d'exécution: (\texttt{ucontext\_t} OU \texttt{jmp\_buf})
    \item Pile associée au thread (8 Ko)
    \item État (\texttt{READY/RUNNING/BLOCKED/TERMINATED})
    \item Valeur de retour (\texttt{void* retval})
    \item Thread en attente (\texttt{joiningThread})
    \item Liaison (\texttt{TAILQ\_ENTRY}) pour ordonnancement
\end{itemize}
Nous avons retenu les \textit{TAILQ} de la bibliothéque \textit{OpenBSD} car ils permettent un ordonnancement FIFO efficace et de supprimer n'importe quel thread en $O(1)$, ce qui nous évitera des parcours linéaires qui vont affecter les performances.

\subsection{Outils de Compilation et Tests}
Le \texttt{Makefile} supporte les cibles suivantes :
\begin{itemize}
    \item \texttt{make} : compile la bibliothèque + tests
    \item \texttt{make check} : exécution de l'ensemble des tests
    \item \texttt{make valgrind} : analyse les fuites mémoire (0 leaks confirmés)
    \item \texttt{make pthreads} : utilise l'implémentation de la libc pour les comparaisons
    \item \texttt{make install} : installation arborescence \texttt{install/}
    \item \texttt{make graphs} : génération des graphes en faisant varier le nombre de threads
\end{itemize}

\subsection{Mutex}
Nous avons opté pour une implémentation simple mais efficace pour les mutex:
\begin{itemize}
    \item Structure : un entier (1 ou 0) + file d'attente \texttt{TAILQ} qui contient l'ensemble des threads qui attendent le déverouillage de ce mutex.
    \item Fonctions : \texttt{init()}, \texttt{destroy()}, \texttt{lock()}, \texttt{unlock()}
\end{itemize}


\subsection{Ordonnancement Coopératif}
\begin{itemize}
    \item \texttt{readyQueue} (\texttt{TAILQ}) : threads prêts à s'exécuter (FIFO)
    \item \texttt{blockedQueue} (\texttt{TAILQ}) : threads en attente (\texttt{join})
    \item \texttt{currentThread} : variable global pointant vers le thread courant
\end{itemize}

